\section{서론}
\label{sec:intro}

종단간(end-to-end, E2E) 자율주행은 센서 관측에서 미래 이동 경로나 제어 의도를 직접 예측하며, 최근 비전--언어 주행 모델은 행동 이유를 설명하는 Chain-of-Causation(CoC)을 함께 생성한다 \cite{hu2023uniad,jiang2023vad,sima2024drivelm,tian2024drivevlm,wang2025alpamayo}. 문제는 같은 장면의 CoC들이 개별적으로는 자연스러워도 시간순으로 함께 읽으면 양립하지 않을 수 있다는 점이다. 예를 들어 ``보행자가 통과할 때까지 정지한다'' 다음에 통과 증거 없이 ``가속하여 진행한다''가 나오면 정지와 진행 요구가 동시에 남는다. 반대로 종료 조건을 만족했는데도 정지 요구가 남으면 정상 진행을 모순으로 오판한다. 이러한 불일치는 학습 자료의 행동 요구를 모호하게 하고 같은 자료를 다시 검사했을 때 일관된 결과를 얻기 어렵게 한다. 하지만 현재 사건만 보는 검사는 이전 사건에서 시작된 요구를 버리므로 이런 오류를 찾을 수 없다.

본 연구는 이 문제를 해결하기 위해 각 CoC를 따로 검사하지 않고, 이전 문장의 요구를 다음 문장까지 기억하며 시간순으로 검사한다. 핵심 발상은 한 CoC가 만든 정지ㆍ양보 요구를 다음 사건으로 전달하고, 그 요구가 끝났다는 명시적인 증거가 나타날 때만 지우는 것이다. 그러면 현재 행동을 앞서 남은 요구와 함께 비교하여 두 문장이 동시에 성립 가능한지 판단하고, 성립하지 않는 최초 사건과 그 이유를 찾을 수 있다. 자연 주석 검토는 사람이 실제 주석을 해석할 때 겪는 어려움을 보여 준다. 이와 별도로, 정상 사례에서 한 항목만 의도적으로 바꾼 합성 비교는 이전 요구를 기억하는지에 따라 판정이 어떻게 달라지는지 확인한다.

\begin{figure*}[t]
\centering
\resizebox{0.95\textwidth}{!}{%
\begin{tikzpicture}[
  main/.style={draw, rounded corners=2pt, align=center, minimum height=12mm,
    text width=3.05cm, font=\footnotesize},
  stage/.style={draw, rounded corners=2pt, align=center, minimum height=8mm,
    text width=3.45cm, font=\scriptsize, fill=black!3},
  data/.style={main, fill=observed!10},
  reason/.style={main, fill=derived!10},
  contract/.style={main, fill=assumed!12},
  route/.style={main, fill=neutral!8},
  flow/.style={-{Latex[length=2mm]}, thick, draw=black!80},
  guide/.style={-{Latex[length=1.5mm]}, thin, dashed, draw=black!45}
]

% Main narrative.
\node[data] (log) at (0,0) {Recorded driving data\\ego trajectory};
\node[reason] (llm) at (4.0,0) {Sequential VLM/LLM\\CoC annotations};
\node[contract] (obl) at (8.0,0) {Prior requirement\\end condition and action order};
\node[contract] (monitor) at (12.0,0) {Python checker retaining\\prior requirements\\UPPAAL checker};
\node[route] (route) at (16.0,0) {Separate results\\natural review\\synthetic rule tests\\response/source checks};

\draw[flow] (log) -- coordinate[midway] (m1) (llm);
\draw[flow] (llm) -- coordinate[midway] (m2) (obl);
\draw[flow] (obl) -- coordinate[midway] (m3) (monitor);
\draw[flow] (monitor) -- coordinate[midway] (m4) (route);

% Process interpretation row.
\node[stage] (s1) at (2.0,-1.55) {CoC labels generated\\after the drive};
\node[stage] (s2) at (6.0,-1.55) {Event ordering\\checking-rule construction};
\node[stage] (s3) at (10.0,-1.55) {True/false/unknown evidence\\remaining-requirement update};
\node[stage] (s4) at (14.0,-1.55) {Decision and first\\violating event};

\draw[guide] (m1) -- (s1.north);
\draw[guide] (m2) -- (s2.north);
\draw[guide] (m3) -- (s3.north);
\draw[guide] (m4) -- (s4.north);

% Takeaway.
\node[draw=assumed, rounded corners=2pt, align=center, font=\footnotesize,
  fill=assumed!8, text width=14.0cm] (key) at (8.0,-3.0)
  {Key point: retaining prior requirements exposes cross-event conflicts, but success on designed synthetic cases does not give accuracy on unseen natural CoCs.};

\end{tikzpicture}%
}

\bifigcaption{연속 CoC에서 지속 요구, 종료 조건, 행동 순서와 미확정 증거를 추출하여 이전 요구를 기억하는 검사 규칙으로 변환하는 전체 흐름. 자연 표본 검토와 합성 규칙 시험은 서로 다른 증거로 보고한다.}{Flow for turning sequential CoCs into rules that retain prior requirements, with natural review separated from controlled synthetic tests.}
\label{fig:overview-pipeline}
\end{figure*}

그림~\ref{fig:overview-pipeline}은 이 발상을 자료 정리, 검사 규칙 만들기, 시간순 판정과 결과 구분의 흐름으로 보여 준다. 먼저 각 CoC에서 행동, 계속 지켜야 할 요구와 그 요구가 끝나는 조건을 읽어 낸다. 이를 입력과 통과ㆍ위반 조건이 분명한 검사 규칙으로 만든다. Python 검사기는 아직 남아 있는 요구와 처음 위반한 사건을 기억한다. 같은 규칙을 별도로 구현한 \uppaal{} 모델 검사기는 정해진 시험 사례에서 최종 판정, 오류 종류, 첫 위반 사건과 정보 부족 사유가 Python 결과와 같은지 확인한다. 따라서 Python은 주 검사기이고, \uppaal{}은 다른 프로그램에서도 같은 판정과 위반 기록을 얻는지 확인하는 보조 검사기이다.

검증은 두 종류로 나뉜다. 중심인 \emph{연속 문장 검사}는 이전 CoC가 만든 요구와 그 요구가 끝나는 조건을 다음 사건까지 기억하여 문장들이 함께 성립할 수 있는지 확인한다. 보조인 \emph{반응 시간 검사}는 개별 CoC가 요구한 행동과 기록된 자차 움직임이 정해 둔 시간 안에 맞아떨어지는지 확인하고, 반복 사건과 장면 경계도 따로 다룬다. 문장들이 서로 모순되지 않아도 실제 움직임이 선택한 시간 기준과 맞지 않을 수 있고, 반대로 움직임이 검출되어도 이전 요구가 문장에 남을 수 있다. 따라서 두 결과를 하나의 정확도나 안전 점수로 합치지 않는다.

본 논문의 연구 질문은 다음 하나이다.
\begin{quote}
\textbf{연구 질문:} 이전 요구를 기억하는 검사는 현재 사건만 보는 검사보다 앞선 사건의 정보가 있어야 드러나는 일관성 오류를 더 많이 찾는가? 또한 같은 판정 규칙을 구현한 \uppaal 검사는 Python과 같은 최종 판정, 오류 종류, 최초 위반 사건과 정보 부족 사유를 내는가?
\end{quote}

본 논문은 저장된 연속 CoC를 시간순으로 검사하는 \emph{일관성 검사 프레임워크}를 제안한다. 핵심 아이디어는 CoC를 서로 독립된 설명으로 보지 않고, 앞선 문장에서 시작된 요구가 이후 문장까지 이어질 수 있는 하나의 흐름으로 표현하는 것이다. 이를 위해 각 CoC에서 요구 행동, 요구가 끝나는 조건, 행동 순서와 판단에 필요한 증거를 읽어 공통 검사 규칙으로 만든다.

검사기는 아직 지켜야 하는 요구를 목록으로 유지한다. 새 요구가 나오면 목록에 넣고, 요구가 끝났다는 분명한 증거가 있을 때만 지우며, 증거가 부족하면 임의로 판단하지 않고 그 이유를 남긴다. 그런 다음 현재 행동을 남아 있는 요구 및 행동 순서와 비교하여, 처음 발견한 일관성 오류와 그 이유를 기록한다. 이전 요구를 기억할 때만 구분할 수 있는 오류는 다음 두 가지이다. 첫째, ``보행자가 지나갈 때까지 정지''라는 요구가 아직 끝나지 않았는데 다음 CoC가 진행을 요구하는 문장 사이의 행동 충돌이다. 둘째, 보행자가 지나가 정지 요구가 끝났는데도 검사기가 그 요구를 목록에서 지우지 않아 이후의 정상적인 진행을 오류로 판단하는 처리 오류이다. 즉 첫째는 입력 문장들이 함께 성립할 수 없는 경우이고, 둘째는 끝난 요구를 잘못 남긴 검사 처리 문제이다. 본 논문은 이 둘과 한 사건 안에서 드러나는 두 규칙 위반을 묶어 \emph{일관성 오류}라고 부른다.

프레임워크의 각 역할은 서로 다른 방법으로 확인한다. 자연 자료 검토는 자동 생성 주석을 사람이 얼마나 일관되게 해석할 수 있는지 살피고, 한 항목만 바꾼 합성 비교는 이전 요구를 기억하는 효과를 다른 입력 차이와 분리한다. Python 검사기는 주 판정을 수행하며, \uppaal{}은 새로운 증명 방법이 아니라 같은 규칙과 고정 사례를 별도로 실행하여 판정과 위반 기록이 같은지 확인한다. 반응 시간 설정과 장면 연결 근거에 관한 분석은 중심 문장 검사와 구분하여, 검사 결과가 어떤 조건에서 달라지는지를 보조적으로 설명한다.

핵심 기여는 세 가지이다. 첫째, 연속 CoC에서 계속 지켜야 할 요구, 그 요구가 끝나는 조건, 행동 순서와 참ㆍ거짓ㆍ미상 증거를 구분한다. 둘째, 이전 요구를 기억하는 Python 검사기와 \uppaal 검사 모델로 최종 판정, 일관성 오류의 종류, 최초 위반 사건과 정보 부족 사유를 기록한다. 셋째, 자연 자료 검토와 통제된 합성 규칙 시험을 분리하여 각 결과가 무엇을 보여 주는지 명확히 한다. 본 논문은 자연 오류율, 물리적 차량 안전 또는 학습 성능 향상을 주장하지 않는다.
